\subsection{\texorpdfstring{\(p\) -rings}{p-rings}}

DEFINITION 1. --- A ring C is said to be a p-ring if the ideal pC of C is maximal, and if C is separated and complete for the pC-adic topology.

Let C be a ring; if \(p1_{C}\) is nilpotent and if the ideal pC of C is maximal, C is a p-ring, for the pC-adic topology of C is discrete. In particular, every field of characteristic p is a p-ring.

PROPOSITION 1. --- Let C be a p-ring.

\begin{enumerate}
\def\labelenumi{\alph{enumi})}
\item
  The ring C is local, with maximal ideal pC.
\item
  Suppose that \(p1_{\mathbb{C}}\) is nilpotent. Let \(d\) be the smallest positive integer such that \(p^{d}1_{\mathbb{C}} = 0\) . The ideals of \(\mathbb{C}\) are of the form \(p^{k}\mathbb{C}\) with \(0 \leqslant k \leqslant d\) and one has \(p^{k}\mathbb{C} \neq p^{l}\mathbb{C}\) when \(k\) and \(l\) are two distinct integers satisfying \(0 \leqslant k \leqslant d, 0 \leqslant l \leqslant d\) . The \(\mathbb{C}\) -module \(\mathbb{C}\) is of length \(d\) .
\item
  Suppose that \(p1_{\mathbb{C}}\) is not nilpotent. Then C is a discrete valuation ring whose residue field is of characteristic \(p\) , and whose field of fractions is of characteristic 0. The ideals of the form \(p^{n}\mathbf{C}\) , with \(n \in \mathbf{N}\) , are pairwise distinct; they form all the nonzero ideals of C. The C-module C is not of finite length.
\end{enumerate}

The assertion \(a)\) follows from prop. 19 of III, § 2, no. 13.

We have \(\bigcap_{n\geqslant 0}p^{n}\mathbf{C} = \{0\}\) by hypothesis. Let \(x\neq 0\) in \(\mathbf{C}\) ; there exists an integer \(n\geqslant 0\) such that \(x\in p^n\mathbf{C}, x\notin p^{n + 1}\mathbf{C}\) ; there therefore exists an element \(y\) of \(\mathbf{C}\) such that \(x = p^n y\) ; since \(y\) does not belong to \(p\mathbf{C}\) , \(y\) is invertible.

Suppose that \(p1_{\mathbb{C}}\) is not nilpotent. If \(x\) and \(x'\) are two nonzero elements of C, there exist two integers \(n \geqslant 0\) , \(n' \geqslant 0\) and two invertible elements \(y\) , \(y'\) of C such that \(x = p^n y\) , \(x' = p^{n'} y'\) . Then one has \(xx' = p^{n+n'} yy' \neq 0\) , hence C is an integral domain. Since C is a local ring, but is not a field, and the maximal ideal \(m_{\mathbb{C}} = p\mathbb{C}\) of C is principal, C is a discrete valuation ring (VI, § 3, no. 6, prop. 9). The nonzero ideals of C are then of the form \(p^n\mathbb{C}\) by loc. cit., prop. 8, and are pairwise distinct. In particular, the ring C is not Artinian, hence the C-module C is not of finite length. The residue field C/pC of C is of characteristic \(p\) . Let \(q\) be the characteristic of the field of fractions of C. We have \(p1_{\mathbb{C}} \neq 0\) , whence \(p \neq q\) . Moreover, if \(q\) were nonzero, we would have \(q1_{\mathbb{C}} = 0\) , hence C/pC would be of characteristic \(q \neq p\) , which is absurd. This proves \(c\) ).

Suppose that \(p1_{C}\) is nilpotent. Let \(d\) be the smallest positive integer such that \(p^{d}1_{C} = 0\) . We have a sequence of ideals

\[
\mathrm{C} \supset p \mathrm{C} \supset p ^ {2} \mathrm{C} \supset \dots \supset p ^ {d - 1} \mathrm{C} \supset p ^ {d} \mathrm{C} = \{0 \}.\tag{E}
\]

If \(k\) is an integer such that \(0 \leqslant k < d\) and \(p^k\mathbf{C} = p^{k+1}\mathbf{C}\) , then we deduce

\[
p ^ {d - k - 1} p ^ {k} \mathrm{C} = p ^ {d - k - 1} p ^ {k + 1} \mathrm{C} = \{0 \}
\]

contrary to the hypothesis \(p^{d-1}1_{\mathbf{C}} \neq 0\) . Hence the elements of the sequence (E) are pairwise distinct. Let \(\mathfrak{a}\) be an ideal of C and let \(k\) be the smallest positive integer such that \(\mathfrak{a} \supset p^k\mathbf{C}\) . Let \(x\) be a nonzero element of \(\mathfrak{a}\) ; we have seen that \(x\) is of the form \(p^m u\) with \(m \geqslant 0\) and \(u\) is invertible in C. We therefore have \(p^m\mathbf{C} \subset \mathfrak{a}\) , whence \(m \geqslant k\) , and finally \(x \in p^k\mathbf{C}\) . In conclusion, we have \(\mathfrak{a} = p^k\mathbf{C}\) . The sequence (E) is then a Jordan-Hölder sequence of the C-module C, which is of length \(d\) .

COROLLARY 1. --- If the p-ring C is an integral domain, it is a discrete valuation ring, or a field of characteristic p.

Suppose C is integral. If \(p1_{\mathbf{C}}\) is nilpotent, we have \(p1_{\mathbf{C}} = 0\) , and \(\{0\}\) is a maximal ideal of C, hence C is a field of characteristic \(p\) . If \(p1_{\mathbf{C}}\) is not nilpotent, then C is a discrete valuation ring by prop. 1, c).

COROLLARY 2. --- Let C be a p-ring and a an ideal of C distinct from C. The ring C/a is a p-ring.

We may assume \(\mathfrak{a} \neq \{0\}\) . There then exists an integer \(i \geqslant 1\) such that \(\mathfrak{a} = p^i\mathrm{C}\) ; the ideal \(p\mathrm{C} / \mathfrak{a}\) of \(\mathrm{C} / \mathfrak{a}\) is maximal and we have \(p^i1_{\mathrm{C} / \mathfrak{a}} = 0\) , hence \(\mathrm{C} / \mathfrak{a}\) is a \(p\) -ring.

Let C be a \(p\) -ring. We call the length of C, and denote by \(l(\mathbf{C})\) , the least upper bound in \(\overline{\mathbf{R}}\) of the set of integers \(n \geqslant 1\) such that \(p^{n-1}1_{\mathbf{C}} \neq 0\) . When \(l(\mathbf{C})\) is finite, it is the length of the C-module C, and when \(l(\mathbf{C})\) is equal to \(+\infty\) , the C-module C is not of finite length (prop. 1).

Examples. --- 1) For every integer \(n \geqslant 1\) , the ring \(\mathbf{Z}/p^{n}\mathbf{Z}\) is a \(p\) -ring of length \(n\) . The ring \(\mathbf{Z}_{p}\) of \(p\) -adic integers is a \(p\) -ring of infinite length.

\begin{enumerate}
\def\labelenumi{\arabic{enumi})}
\setcounter{enumi}{1}
\tightlist
\item
  Let K be a perfect field of characteristic p. By prop. 8 of § 1, no. 8, the ring W(K) of Witt vectors is a p-ring of infinite length. The map \((a_{n})_{n\in\mathbf{N}} \mapsto a_{0}\) induces by passage to the quotient an isomorphism of W(K)/pW(K) onto the field K (loc. cit., prop. 7). For every integer \(n \geqslant 1\) , the ring
\end{enumerate}

\[
\mathrm{W} _ {n} (\mathrm{K}) = \mathrm{W} (\mathrm{K}) / p ^ {n} \mathrm{W} (\mathrm{K})
\]

is a p-ring of length n.

PROPOSITION 2. --- Let C and C' be two p-rings and let \(u\) be a homomorphism from C to C'. Let \(v\) be the homomorphism from \(\kappa_{\mathbf{C}} = \mathbf{C} / p\mathbf{C}\) to \(\kappa_{\mathbf{C}'} = \mathbf{C}' / p\mathbf{C}'\) induced by u by passage to the quotients.

\begin{enumerate}
\def\labelenumi{\alph{enumi})}
\item
  We have \(l(\mathbf{C}) \geqslant l(\mathbf{C}')\) and u is injective if and only if \(l(\mathbf{C}) = l(\mathbf{C}')\) .
\item
  For u to be surjective, it is necessary and sufficient that v be an isomorphism.
\item
  For u to be an isomorphism, it is necessary and sufficient that v be an isomorphism and that we have \(l(\mathbf{C}) = l(\mathbf{C}')\) .
\end{enumerate}

Let \(n \geqslant 1\) be an integer. One has \(u(p^{n-1}1_{\mathbf{C}}) = p^{n-1}1_{\mathbf{C}'}\) , hence the relation \(p^{n-1}1_{\mathbf{C}'} \neq 0\) implies \(p^{n-1}1_{\mathbf{C}} \neq 0\) and is equivalent to it if \(u\) is injective. We therefore have \(l(\mathbf{C}') \leqslant l(\mathbf{C})\) with equality if \(u\) is injective. If \(u\) is not injective, there exists an integer \(i < l(\mathbf{C})\) such that the kernel of \(u\) is the ideal \(p^i\mathbf{C}\) of \(\mathbf{C}\) ; we then have \(p^i1_{\mathbf{C}'} = 0\) , whence \(l(\mathbf{C}') \leqslant i\) . This proves \(a\) ).

Since \(\kappa_{\mathrm{C}}\) and \(\kappa_{\mathrm{C}'}\) are fields, the homomorphism \(v\) is injective. If \(u\) is surjective, so is \(v\) , which is therefore an isomorphism. Conversely, suppose \(v\) is surjective. Then for every integer \(n \geqslant 0\) , the map \(v_{n}: p^{n}\mathrm{C}/p^{n+1}\mathrm{C} \to p^{n}\mathrm{C}'/p^{n+1}\mathrm{C}'\) derived from \(u\) is surjective. Since C is complete for the \(p\mathrm{C}\)-adic filtration and \(\mathbf{C}'\) separated for the \(p\mathbf{C}'\)-adic filtration, \(u\) is surjective by Cor. 2 of Th. 1 of III, § 2, No. 8. This proves \(b\) ).

Finally, c) follows from a) and b).

PROPOSITION 3. --- Let \((\mathbf{C}_{n}, \pi_{n,m})\) be a projective system of rings relative to the index set N. We assume that \(\mathbf{C}_{n}\) is a p-ring for all \(n \in \mathbf{N}\) and that the homomorphisms \(\pi_{n,m}\) are surjective. Then \(C = \varprojlim C_{n}\) is a p-ring, and for all \(n \in \mathbf{N}\) the canonical homomorphism \(\pi_{n}: C \to C_{n}\) is surjective and induces an isomorphism of \(\kappa_{C}\) onto \(\kappa_{C_{n}}\) .

As the maps \(\pi_{n,m}\) are surjective, the same holds for the maps \(\pi_n\) (E, III, p. 58, prop. 5). Let us show that C is a \(p\) -ring. Let \(d_n\) be the length of \(\mathbf{C}_n\) . By prop. 2, a), the sequence of elements \(d_n\) of \(\mathbf{N} \cup \{+\infty\}\) is increasing; if it is stationary, there exists an integer \(n_0\) such that \(\pi_{n,m}\) is an isomorphism of \(\mathbf{C}_m\) onto \(\mathbf{C}_n\) when \(n_0 \leqslant n \leqslant m\) , so that C, which is isomorphic to \(\mathbf{C}_{n_0}\) , is a \(p\) -ring.

It therefore suffices to consider the case where each \(d_{n}\) is finite, and where the sequence \((d_{n})\) tends to \(+\infty\) . Endow the ring C with the trivial filtration (III, § 2, n° 1, example 5). For \(n\in\mathbf{N}\) , let \(\mathbf{I}_{n}\) be the kernel of \(\pi_{n}\) ; set \(\mathbf{I}_{n}=\mathbf{C}\) if \(n<0\) . Let E denote the C-module C equipped with the filtration \((\mathbf{I}_{n})_{n\in\mathbf{Z}}\) . It is separated and complete, because the topology \(\mathcal{G}\) defined by the filtration \((\mathbf{I}_{n})_{n\in\mathbf{Z}}\) is the projective limit topology of the discrete topologies on the \(\mathbf{C}_{n}\) .

Let \(k\) be an integer \(\geqslant 1\) . We have \(p^{k}\mathrm{C} \subset \varprojlim (p^{k}\mathrm{C}_{n})\) (E, III, p. 55, formula (9)). Conversely, if \(x = (x_{n})_{n \in \mathbb{N}} \in \varprojlim (p^{k}\mathrm{C}_{n})\) and if one sets \(\mathrm{X}_{n} = \{y \in \mathrm{C}|\pi_{n}(p^{k}y) = x_{n}\}\) , the sequence \((\mathrm{X}_{n})_{n \in \mathbb{N}}\) is a decreasing sequence of nonempty closed affine subsets of E. Since \(\mathrm{E / I}_{n}\) is an artinian C-module, the intersection of the \(\mathrm{X}_{n}\) is nonempty (III, § 2, n° 7, prop. 7); for every \(z \in \bigcap_{n \in \mathbb{N}} \mathrm{X}_{n}\) we have \(p^{k}z = x\) . We have therefore proved that one has \(p^{k}\mathrm{C} = \varprojlim p^{k}\mathrm{C}_{n}\) .

for every integer \(k \geqslant 1\) . In particular, the ideal \(p^{k}\text{C}\) of C is closed for the topology \(\mathcal{G}\) . On C, the \(p\) -adic topology is finer than the topology \(\mathcal{G}\) because one has \(p^{d_{n}}\text{C} \subset \text{I}_{n}\) . It then follows from TG, III, p. 26, cor. 1 to prop. 10, that C is separated and complete for the \(p\text{C}\) -adic topology. Moreover, one has \(p\text{C} = \varprojlim p\text{C}_{n} = \pi_{0}^{-1}(p\text{C}_{0})\) and hence the surjective homomorphism of C/pC into \(\text{C}_{0}/p\text{C}_{0}\) deduced from \(\pi_{0}\) is an isomorphism. This shows that the ideal \(p\text{C}\) of C is maximal and consequently that C is a \(p\) -ring. The last assertion of prop. 3 follows from prop. 2, b).

